\chapter{Introduction}

\section{Problem statement}

Evolutionary design studies the automatic generation and improvement of artefacts whose structure is evaluated by a task-dependent fitness function. In this thesis, the considered artefacts are virtual three-dimensional organisms represented as Framsticks F1 genotypes and evaluated with respect to the \texttt{vertpos} criterion. The genotype is not a fixed-length vector in which every coordinate has an independent and stable meaning. Instead, an F1 genotype is a variable-length symbolic expression with a branching, tree-like interpretation. The F1 representation combines sequential elements, branches, cores, and modifiers, which together define the syntactic structure from which a phenotype is constructed. Consequently, the same local symbol or fragment may have a different effect depending on its position and surrounding genotype context.

This representation creates a difficult optimization problem. Search is performed in a discrete and highly constrained space, because candidate strings must remain syntactically valid and must decode into valid Framsticks structures accepted for evaluation before their fitness can be assessed. Validity and semantics are therefore not properties of isolated characters only, but of complete structured expressions. A local modification that is syntactically acceptable in one part of a genotype may be invalid, neutral, or harmful in another part. This context dependence limits the effectiveness of naive edits and makes the definition of useful neighbourhoods non-trivial.

The standard evolutionary approach in Framsticks provides a natural and strong baseline for this type of problem. However, the thesis investigates whether additional representations or modification principles can improve the search process or at least offer competitive alternatives. The first approach, Compatible Substitutions Optimization (CoSO), addresses the problem directly in the discrete genotype representation by considering whether a proposed substitution is compatible with the context in which it is inserted. In this thesis, CoSO is treated as a context-aware family of genotype-editing methods.

The second approach, Latent Space Optimization (LSO), attempts to transform the same structured genotype search problem into optimization over a fixed-dimensional continuous representation learned by autoencoders. In this setting, a model encodes an F1 genotype into a latent vector, search operators modify vectors in the latent space, and the decoder maps the resulting vectors back to candidate genotypes that can be evaluated in Framsticks. This strategy is promising because it may introduce smoother neighbourhood relations than raw symbolic edits, but it also introduces new risks: decoded candidates may be invalid, the learned geometry may not correspond to fitness-improving changes, and good reconstruction quality does not by itself guarantee good optimization performance.

The problem addressed in this thesis is therefore the comparative assessment of two ways of handling context-dependent, variable-length genotypes in evolutionary design. CoSO represents an explicit, discrete, context-aware strategy, whereas LSO represents an implicit, learned, continuous strategy. The central question is not whether either approach universally replaces the standard Framsticks evolutionary algorithm, but whether and when these approaches can provide competitive search behaviour relative to this baseline for F1 genotypes optimized for \texttt{vertpos}. Detailed evaluation criteria and experimental comparisons are defined in the later chapters.

\section{Motivation}

Evolutionary design is motivated by the possibility of producing complex structures automatically, without requiring a designer to specify every morphological detail in advance. Systems such as Framsticks are particularly suitable for this type of research because they combine an explicit genotype representation, a simulation-based phenotype evaluation, and open-ended structural variation. At the same time, these properties make optimization difficult. A useful method should not only discover high-fitness organisms, but also preserve enough syntactic and structural coherence for the search process to remain productive over many generations.

The F1 representation highlights a broader challenge that appears in many generative and artificial-life settings: candidate solutions are neither simple vectors nor independent collections of parameters. They are structured objects whose local fragments interact through grammar, branching, and phenotype construction. This makes the design of variation operators especially important. Random local changes may explore the space widely, but they can also destroy useful substructures or produce candidates that cannot be decoded into valid structures for simulation-based evaluation. Therefore, there is practical motivation to study methods that use more information about the representation than purely position-independent mutation.

Compatible Substitutions Optimization (CoSO) is motivated by this need for explicit context awareness. Instead of treating a substitution as an isolated edit, CoSO considers whether the inserted fragment is compatible with the surrounding genotype context. This perspective is methodologically useful because it remains close to the original discrete representation and can be interpreted in terms of concrete genotype fragments. In evolutionary design, such interpretability is valuable: it can support explanations of why certain modifications may be more or less suitable in a given context. However, the usefulness of this idea must be assessed empirically, because compatibility with the local context does not automatically imply improvement in a simulation-based fitness criterion such as \texttt{vertpos}.

Latent Space Optimization (LSO) is motivated by a complementary observation. Instead of defining all useful neighbourhood relations manually in the symbolic genotype space, an autoencoder may learn a continuous representation in which related genotypes are mapped to nearby vectors. If this representation captures relevant regularities of F1 genotypes, standard continuous or population-based search operators can be applied in the latent space and then decoded back to candidate genotypes. Models that respect syntactic and structural regularities are especially relevant in this setting, but the methodological details belong to the later background and methods chapters. Nevertheless, LSO also requires careful evaluation: reconstruction quality, syntactic validity, and optimization performance are related but distinct properties.

The motivation for comparing CoSO and LSO is therefore not that one approach is expected to universally replace the other or to dominate the standard Framsticks evolutionary algorithm. Rather, they represent two different ways of addressing the same structural difficulty. CoSO makes context dependence explicit in the discrete search space, whereas LSO attempts to learn a continuous search space in which some of this context dependence may be embedded implicitly. Studying both approaches under a common case study can clarify which benefits come from explicit compatibility information, which come from learned representations, and which limitations remain inherent to the optimization of variable-length F1 genotypes.

\section{Aim and scope of the thesis}

The main aim of this thesis is to compare two representation-aware approaches to optimization in evolutionary design: Compatible Substitutions Optimization (CoSO) and Latent Space Optimization (LSO). Both approaches address the same underlying difficulty of the Framsticks F1 representation, namely that useful modifications of a genotype depend on syntax, branching structure, and surrounding context. The comparison is not intended to prove that either approach universally replaces the standard Framsticks evolutionary algorithm. Instead, the thesis aims to determine how explicit context-aware substitutions and learned continuous representations behave when they are applied to the same class of variable-length genotypes and evaluated with the same objective.

The empirical scope of the work is limited to Framsticks F1 genotypes optimized for the \texttt{vertpos} fitness criterion. This restriction makes it possible to define a common case study in which all candidate solutions are ultimately evaluated by the same simulation-based criterion. Within this scope, CoSO represents the discrete and explicit direction of the study, while LSO represents the learned continuous-representation direction.

The LSO part of the thesis covers preparing F1 genotype data for autoencoder training, representing genotypes with character-based or grammar-aware encodings, training selected autoencoder architectures, decoding latent vectors back into F1 genotypes, and applying optimization methods in the learned latent space. The autoencoder-based investigation focuses on whether the learned representation can support productive search, not only on whether it reconstructs training examples accurately.

The thesis also includes a baseline comparison with Framsticks EA. This baseline is important because it represents the native evolutionary search procedure for the considered environment and provides a practical reference point for both CoSO and LSO. The intended outcome of the study is therefore a comparative assessment of solution quality, validity, stability, and dependence on starting conditions, rather than a single claim that one method is always superior.

\section{Proposed research approach}

The research approach follows a common experimental case study centered on F1 genotypes and the \texttt{vertpos} objective. The first step is to establish the representation and evaluation setting: genotypes are treated as structured symbolic expressions, and candidate solutions are evaluated through Framsticks. This provides a shared basis for comparing methods that operate in different search spaces.

For CoSO, the thesis considers optimization through context-aware substitutions in the original discrete genotype representation. This approach is analysed as an explicit attempt to use information about the compatibility of genotype fragments with their surrounding context.

For LSO, the research approach consists of learning continuous representations of F1 genotypes with autoencoders and then applying search operators to vectors in the latent space. A candidate latent vector is decoded into an F1 genotype and evaluated with the same Framsticks fitness criterion as the baseline. This makes it possible to compare latent-space search with direct genotype-level evolution while preserving a common final evaluation procedure. The implemented LSO component considers several families of autoencoder models and several continuous or population-based optimization strategies, but their technical details are introduced only in the methods and experimental chapters.

The comparison is designed to separate three related questions. First, it examines whether generated or modified genotypes remain valid and evaluable. Second, it assesses whether the search process improves the fitness of starting genotypes. Third, it compares the stability of methods across different starting conditions. This structure reflects the fact that a representation can be syntactically convenient without necessarily producing better optimization results, and that a high-performing single run is not sufficient evidence of a stable optimization method.

\section{Thesis structure}

The remainder of the thesis is organized as follows. Chapter~2 presents the literature and background required for the study. It introduces evolutionary design with variable-length genotypes, the Framsticks environment and the F1 representation, context-aware genotype modification, Compatible Substitutions Optimization, autoencoders, variational autoencoders, and the general idea of Latent Space Optimization.

Chapter~3 describes the methods and implemented components. It covers the dataset and fitness evaluation procedure, the grammar-based handling of F1 genotypes, the CoSO component, the autoencoder representations, latent-space optimization methods, baseline evolutionary methods, and the software architecture used to run the experiments.

Chapter~4 defines the experimental design. It formulates the research questions, identifies the compared methods, describes the benchmark setup and evaluation metrics, and explains the experimental protocol and reproducibility considerations. Chapter~5 reports the experimental results, including reconstruction and validity analyses for autoencoders, CoSO results, latent-space optimization results, comparisons with Framsticks EA, a direct comparison between CoSO and Latent Space Optimization, and analyses conditioned on the initial fitness range.

Chapter~6 discusses the results and concludes the thesis. It summarizes the main findings, analyses the strengths and limitations of CoSO and LSO, compares the two approaches, discusses threats to validity, and outlines possible directions for future work. Supplementary material is provided in the appendix where appropriate.
